\apendice{Documentación de usuario}

\section{Introducción}

Este apéndice recoge la documentación orientada al usuario final de la aplicación desarrollada en este Trabajo de Fin de Grado. A diferencia de la documentación técnica de programación, esta sección no describe la instalación interna del sistema, la configuración de Docker, la base de datos ni las dependencias del proyecto.

Desde el punto de vista del usuario, la aplicación se utiliza como una plataforma web. Por tanto, una vez desplegada por el personal técnico o por el administrador del sistema, el usuario solo necesita acceder a la dirección proporcionada mediante un navegador web.

El objetivo de este apéndice es explicar los requisitos básicos de acceso y dejar preparada la sección posterior del manual de usuario, donde se describirá el uso funcional de la aplicación: registro, inicio de sesión, ejecución de pipelines, consulta de resultados, gestión de compras y funciones de administración.

\section{Requisitos de usuarios}

Para utilizar la aplicación no es necesario instalar Python, MariaDB, Docker, Docker Compose ni ninguna dependencia interna del proyecto. Todos esos elementos forman parte del despliegue técnico y se documentan en el apéndice de documentación técnica de programación.

El usuario únicamente necesita disponer de:

\begin{itemize}
	\item \textbf{Navegador web actualizado}: necesario para acceder a la interfaz de la aplicación.
	\item \textbf{Conexión a Internet}: o acceso a la red en la que esté desplegada la plataforma.
	\item \textbf{Dirección de acceso}: URL proporcionada por el administrador o responsable técnico.
\end{itemize}

En un entorno local de demostración, la aplicación puede estar disponible en:

\begin{bloqueCodigo}[title={url de acceso}]
	http://localhost:5000
\end{bloqueCodigo}

En caso de acceder desde otro equipo de la misma red, deberá utilizarse la dirección IP o URL proporcionada por el administrador del sistema.

En un despliegue real, el acceso se realizaría mediante la URL pública o privada configurada para la plataforma.

\section{Instalación}

Desde el punto de vista del usuario final, no es necesario instalar la aplicación ni preparar ningún componente técnico en su equipo.

La aplicación se utiliza directamente desde un navegador web, accediendo a la dirección proporcionada por el administrador o responsable del sistema. La instalación, configuración y despliegue de la plataforma corresponden al personal técnico y se describen en la documentación técnica de programación.

Por tanto, el usuario solo necesita disponer de un navegador web actualizado, conexión a Internet o acceso a la red donde esté desplegada la aplicación, y una cuenta válida para iniciar sesión.

La aplicación se ejecuta escuchando en \texttt{0.0.0.0:5000} para permitir conexiones desde fuera del contenedor. Sin embargo, cuando se accede desde el propio equipo anfitrión, la URL habitual de acceso es \texttt{http://localhost:5000}.

\section{Manual del usuario}

\subsection{Problemas frecuentes de acceso}

Los problemas más habituales desde el punto de vista del usuario se resumen en la Tabla~\ref{tabla:problemas_acceso_usuario}.

\begin{table}[H]
	\centering
	\small
	\renewcommand{\arraystretch}{1.15}
	\begin{tabularx}{\linewidth}{p{0.35\linewidth}X}
		\toprule
		\textbf{Situación} & \textbf{Posible causa o solución} \\
		\midrule
		La página no carga & Comprobar la conexión a Internet y que la URL introducida sea correcta. \\
		La URL no es válida & Solicitar al administrador la dirección actual de la aplicación. \\
		No se puede iniciar sesión & Revisar usuario y contraseña. Si el problema continúa, contactar con el administrador. \\
		La cuenta aparece como no activa & La cuenta puede estar suspendida, dada de baja o pendiente de revisión. \\
		No aparecen pipelines disponibles & Puede que el usuario no tenga contratado el plan o modelo necesario. \\
		No se pueden descargar resultados & El archivo puede haber caducado o no estar disponible. \\
		\bottomrule
	\end{tabularx}
	\caption[Problemas frecuentes de acceso]{Problemas frecuentes de acceso para el usuario.}
	\label{tabla:problemas_acceso_usuario}
\end{table}
